% !TeX program = xelatex
\documentclass[a4paper,14pt]{extarticle}

% ---------- GOST-like layout ----------
\usepackage[a4paper,left=30mm,right=10mm,top=20mm,bottom=20mm]{geometry}
\usepackage{setspace}
\onehalfspacing
\setlength{\parindent}{1.25cm}
\setlength{\parskip}{0pt}
\usepackage{indentfirst}
\usepackage[protrusion=false,expansion=false]{microtype}

% ---------- Fonts / language ----------
\usepackage{fontspec}
\setmainfont{Times New Roman}
\setmonofont{Courier New}
\usepackage{polyglossia}
\setdefaultlanguage{russian}
\setotherlanguage{english}
\newfontfamily\cyrillicfont{Times New Roman}
\newfontfamily\cyrillicfonttt{Courier New}

% ---------- Math / tables / graphics ----------
\usepackage{amsmath,amssymb}
\usepackage{graphicx}
\usepackage{float}
\usepackage{booktabs}
\usepackage{array}
\usepackage{caption}
\usepackage{tikz}
\usetikzlibrary{arrows.meta,positioning,calc}

\renewcommand{\figurename}{Рисунок}
\renewcommand{\tablename}{Таблица}
\captionsetup{labelsep=endash,justification=centering}

% ---------- Hyperlinks ----------
\usepackage[hidelinks]{hyperref}
\usepackage{xurl}
\urlstyle{tt}
\DeclareRobustCommand{\code}[1]{\nolinkurl{#1}}

% Better line breaking
\pretolerance=1000
\tolerance=2000
\setlength{\emergencystretch}{2em}
\hbadness=10000

% ---------- Listings ----------
\usepackage{xcolor}
\usepackage{listings}
\renewcommand{\lstlistingname}{Листинг}
\lstset{
  basicstyle=\fontsize{10}{12}\ttfamily,
  keywordstyle=\bfseries,
  commentstyle=\itshape,
  stringstyle=\itshape,
  numbers=left,
  numberstyle=\tiny,
  stepnumber=1,
  numbersep=8pt,
  tabsize=2,
  breaklines=true,
  breakatwhitespace=false,
  showstringspaces=false,
  frame=single,
  rulecolor=\color{black},
  xleftmargin=12pt,
  framexleftmargin=12pt
}

\lstdefinelanguage{cmake}{
  morekeywords={cmake_minimum_required,project,set,add_library,add_executable,target_link_libraries,target_compile_definitions,option,if,elseif,else,endif,find_package,include,FetchContent_Declare,FetchContent_MakeAvailable,add_subdirectory},
  sensitive=true,
  morecomment=[l]{\#},
  morestring=[b]"
}
\lstdefinelanguage{json}{
  morestring=[b]",
  showstringspaces=false,
  morecomment=[l]{//}
}

\begin{document}

\pagenumbering{gobble}
\begin{titlepage}
\begin{center}
\large
САНКТ-ПЕТЕРБУРГСКИЙ НАЦИОНАЛЬНЫЙ\\
ИССЛЕДОВАТЕЛЬСКИЙ УНИВЕРСИТЕТ\\
ИНФОРМАЦИОННЫХ ТЕХНОЛОГИЙ, МЕХАНИКИ И ОПТИКИ\\

\vfill
\textbf{ОТЧЁТ}\par
по дисциплине \textbf{«Архитектура игровых движков»}\\[2mm]
\textbf{«Интеграция физики и системы ввода с визуальной отладкой»}\\[10mm]

\vfill
\begin{flushright}
\begin{tabular}{@{}l@{\quad}l@{}}
Студент: & Тереховский Арсений Дмитриевич\\
Группа: & J3300\\
\end{tabular}
\end{flushright}

\vfill
Санкт-Петербург\\
2026\,г.
\end{center}
\end{titlepage}

\newpage
\pagenumbering{arabic}
\setcounter{page}{2}
\tableofcontents
\newpage

\section{Введение}
Практическое занятие №4 посвящено интеграции физической подсистемы и системы ввода в существующий каркас игрового движка, построенный в предыдущих работах. В рамках задания необходимо было расширить ECS-архитектуру компонентами физики, реализовать обновление тел под действием гравитации, добавить обработку коллизий, систему событий столкновений, визуальную отладку коллайдеров и управление 3D-камерой.

В актуальной версии \textit{TurboGameEngine} эта задача решена на базе модулей \code{engine/ecs}, \code{engine/platform}, \code{engine/game} и \code{engine/renderer}. Рендеринг выполняется через \texttt{RenderAdapter} и backend на \texttt{Diligent Engine}, физическая симуляция работает поверх ECS-компонент \texttt{Transform}/\texttt{Rigidbody}/\texttt{Collider}, а пользовательский ввод обрабатывается через \texttt{platform::InputManager} и контроллер свободной камеры.

\section{Постановка задачи}
\subsection{Цель и задачи}
Согласно выданному заданию, необходимо:
\begin{enumerate}
  \item разработать компоненты \texttt{Rigidbody} и \texttt{Collider} с поддержкой AABB и сферы;
  \item реализовать \texttt{PhysicsSystem}, выполняющую применение гравитации, интегрирование движения, детектирование и разрешение коллизий;
  \item создать \texttt{InputManager} для унифицированной работы с клавиатурой и мышью;
  \item добавить систему событий столкновений;
  \item обеспечить debug-отрисовку коллайдеров через \texttt{RenderAdapter};
  \item реализовать управление 3D-камерой с помощью ввода;
  \item продемонстрировать работу физики и ввода на 3D-сцене.
\end{enumerate}

\subsection{Что реализовано в текущем проекте}
Минимальные требования задания не только выполнены, но и расширены:
\begin{itemize}
  \item помимо AABB-AABB реализованы пересечения \texttt{Sphere-Sphere} и \texttt{AABB-Sphere};
  \item система разрешения столкновений использует не только position correction, но и импульсную модель с трением и коэффициентом восстановления;
  \item добавлены sleeping тел и grid broadphase для снижения количества пар коллизий;
  \item debug-визуализация поддерживает как AABB, так и сферические коллайдеры;
  \item демо-сцена содержит управляемую камеру, динамические кубы, пирамиды и шары.
\end{itemize}

\section{Архитектура решения}
\subsection{Место подсистем физики и ввода в движке}
Физика и ввод встроены в уже существующий игровой цикл следующим образом: \texttt{Application} принимает платформенные события, передаёт их в \texttt{InputManager} и рендерер, затем активное состояние \texttt{GameplayState} использует накопленное состояние ввода для управления камерой и запуском демонстрации. На каждом шаге обновления вызывается \texttt{PhysicsSystem}, а затем сцена отрисовывается обычной \texttt{RenderSystem} и, при необходимости, \texttt{DebugRenderSystem}.

\begin{figure}[H]
\centering
\resizebox{\linewidth}{!}{%
\begin{tikzpicture}[
  node distance=10mm and 10mm,
  box/.style={draw, rounded corners=2pt, align=center, minimum height=10mm, text width=37mm},
  wide/.style={draw, rounded corners=2pt, align=center, minimum height=10mm, text width=46mm},
  arr/.style={-Latex, thick}
]
\node[wide] (app) {\textbf{Application}\\poll events / state stack / frame loop};
\node[box, below left=12mm and 2mm of app] (input) {\textbf{InputManager}\\клавиатура, мышь, scroll, capture};
\node[box, below right=12mm and 2mm of app] (state) {\textbf{GameplayState}\\камера, запуск сцены, ImGui};

\node[box, below=12mm of state, xshift=-24mm] (physics) {\textbf{PhysicsSystem}\\gravity, collision, impulses};
\node[box, right=10mm of physics] (events) {\textbf{EventBus}\\CollisionEnter/Stay/Exit};
\node[box, below=12mm of physics] (render) {\textbf{RenderSystem}\\основная отрисовка};
\node[box, right=10mm of render] (debug) {\textbf{DebugRenderSystem}\\AABB / sphere rings};
\node[wide, below=12mm of debug, xshift=-22mm] (adapter) {\textbf{RenderAdapter / Renderer}\\drawMesh, drawDebugAabb, Diligent backend};

\draw[arr] (app) -- (input);
\draw[arr] (app) -- (state);
\draw[arr] (state) -- (physics);
\draw[arr] (physics) -- (events);
\draw[arr] (state) |- (render);
\draw[arr] (state) |- (debug);
\draw[arr] (render) -- (adapter);
\draw[arr] (debug) -- (adapter);
\draw[arr] (input.east) -| (state.west);
\end{tikzpicture}
}
\caption{Поток данных между подсистемами ввода, физики и рендера}
\label{fig:pz4-arch}
\end{figure}

\subsection{Диаграмма ключевых компонентов и систем}
На уровне ECS данные представлены простыми компонентами, а логика сосредоточена в системах. В текущей работе центральными элементами стали \texttt{Rigidbody}, \texttt{Collider}, \texttt{PhysicsSystem}, \texttt{DebugRenderSystem} и \texttt{FreeCameraController}.

\begin{figure}[H]
\centering
\resizebox{\linewidth}{!}{%
\begin{tikzpicture}[
  node distance=10mm and 12mm,
  cls/.style={draw, rounded corners=2pt, align=left, minimum height=12mm, text width=46mm},
  rel/.style={-Latex, thick}
]
\node[cls] (tr) {\textbf{Transform}\\position\\rotationEulerRadians\\scale};
\node[cls, right=of tr] (rb) {\textbf{Rigidbody}\\velocity, acceleration\\mass, inverseMass\\damping, restitution, friction\\sleeping flags};
\node[cls, right=of rb] (col) {\textbf{Collider}\\type: AABB / Sphere\\offset\\halfExtents / radius};

\node[cls, below=of tr] (cam) {\textbf{Camera}\\FOV, near/far plane\\active flag};
\node[cls, below=of rb] (phys) {\textbf{PhysicsSystem}\\integrateBody()\\buildBroadphasePairs()\\resolveCollision()\\event dispatch};
\node[cls, below=of col] (dbg) {\textbf{DebugRenderSystem}\\draw AABB\\draw sphere rings};

\node[cls, below=of phys] (inp) {\textbf{InputManager}\\raw/captured states\\mouse delta\\scroll delta};
\node[cls, left=of inp] (fcc) {\textbf{FreeCameraController}\\mouse look\\WASDQE fly movement};
\node[cls, right=of inp] (eb) {\textbf{EventBus}\\CollisionEnter\\CollisionStay\\CollisionExit};

\draw[rel] (rb) -- (phys);
\draw[rel] (col) -- (phys);
\draw[rel] (tr) -- (phys);
\draw[rel] (col) -- (dbg);
\draw[rel] (tr) -- (dbg);
\draw[rel] (inp) -- (fcc);
\draw[rel] (cam) -- (fcc);
\draw[rel] (phys) -- (eb);
\end{tikzpicture}
}
\caption{Ключевые компоненты и системы ПЗ4}
\label{fig:pz4-class}
\end{figure}

\section{Реализация ключевых модулей}
\subsection{Компоненты Rigidbody и Collider}
Физические свойства объекта собраны в компоненте \texttt{Rigidbody}. В отличие от минимального варианта из методички, текущая реализация содержит не только линейную скорость и массу, но и угловую скорость, накопленные силы, демпфирование, коэффициенты трения и восстановления, а также поддержку ``засыпания'' тел. Компонент \texttt{Collider} поддерживает два типа коллайдеров --- \texttt{Aabb} и \texttt{Sphere}.

\begin{lstlisting}[language=C++,caption={Компоненты \code{Rigidbody} и \code{Collider}}]
struct Rigidbody final {
    glm::vec3 velocity{0.0F};
    glm::vec3 acceleration{0.0F};
    glm::vec3 accumulatedForce{0.0F};
    glm::vec3 angularVelocity{0.0F};
    glm::vec3 accumulatedTorque{0.0F};
    float mass{1.0F};
    float inverseMass{1.0F};
    float linearDamping{0.08F};
    float angularDamping{0.6F};
    float restitution{0.06F};
    float friction{0.7F};
    float sleepTimer{0.0F};
    bool useGravity{true};
    bool isStatic{false};
    bool isKinematic{false};
    bool canSleep{true};
    bool isSleeping{false};
};

enum class ColliderType : std::uint8_t {
    Aabb = 0,
    Sphere
};

struct Collider final {
    ColliderType type{ColliderType::Aabb};
    glm::vec3 offset{0.0F};
    AabbCollider aabb{};
    SphereCollider sphere{};
    bool enabled{true};
};
\end{lstlisting}

\subsection{Преобразование коллайдера в мировое пространство}
Поскольку ECS хранит коллайдер в локальных координатах сущности, перед детекцией столкновений необходимо получить мировой объём. Для AABB используется абсолютная матрица поворота-масштаба, чтобы получить world-space half-extents; для сферы берётся максимальный масштаб по осям.

\begin{lstlisting}[language=C++,caption={Перевод локального \code{Collider} в мировой объём}]
bool computeWorldColliderShape(
    const World& world,
    const EntityId entity,
    const Collider& collider,
    WorldColliderShape& outShape) {
    const glm::mat4 worldMatrix = computeWorldMatrix(world, entity);
    const glm::vec3 worldCenter = transformPoint(worldMatrix, collider.offset);

    switch (collider.type) {
        case ColliderType::Aabb: {
            const glm::mat3 absBasis = absoluteRotationScale(worldMatrix);
            const glm::vec3 worldHalfExtents = absBasis * collider.aabb.halfExtents;
            outShape.aabb.min = worldCenter - worldHalfExtents;
            outShape.aabb.max = worldCenter + worldHalfExtents;
            return true;
        }
        case ColliderType::Sphere: {
            const glm::vec3 worldScale = extractWorldScale(worldMatrix);
            const float scale = std::max(worldScale.x, std::max(worldScale.y, worldScale.z));
            outShape.sphere.center = worldCenter;
            outShape.sphere.radius = collider.sphere.radius * scale;
            return outShape.sphere.radius > 0.0F;
        }
    }

    return false;
}
\end{lstlisting}

\subsection{Обнаружение коллизий}
Задание требовало минимум AABB-AABB, однако в проекте были добавлены также \texttt{Sphere-Sphere} и \texttt{AABB-Sphere}. Для коробок используется классическая проверка overlap по осям с вычислением минимальной глубины проникновения и нормали разделения.

\begin{lstlisting}[language=C++,caption={Проверка пересечения двух мировых AABB}]
bool intersectAabbAabb(
    EntityId entityA,
    const WorldAabb& aabbA,
    EntityId entityB,
    const WorldAabb& aabbB,
    CollisionContact& outContact) {
    const glm::vec3 overlap = glm::min(aabbA.max, aabbB.max) -
                              glm::max(aabbA.min, aabbB.min);
    if (overlap.x <= 0.0F || overlap.y <= 0.0F || overlap.z <= 0.0F) {
        return false;
    }

    const glm::vec3 centerDelta = aabbB.center() - aabbA.center();
    outContact = CollisionContact{};
    outContact.entityA = entityA;
    outContact.entityB = entityB;
    outContact.contactPoint =
        (glm::max(aabbA.min, aabbB.min) + glm::min(aabbA.max, aabbB.max)) * 0.5F;

    if (overlap.x <= overlap.y && overlap.x <= overlap.z) {
        outContact.penetrationDepth = overlap.x;
        outContact.normal = glm::vec3(centerDelta.x < 0.0F ? -1.0F : 1.0F, 0.0F, 0.0F);
        return true;
    }
    if (overlap.y <= overlap.z) {
        outContact.penetrationDepth = overlap.y;
        outContact.normal = glm::vec3(0.0F, centerDelta.y < 0.0F ? -1.0F : 1.0F, 0.0F);
        return true;
    }

    outContact.penetrationDepth = overlap.z;
    outContact.normal = glm::vec3(0.0F, 0.0F, centerDelta.z < 0.0F ? -1.0F : 1.0F);
    return true;
}
\end{lstlisting}

Такой контакт затем передаётся в фазу разрешения столкновения. Наличие \texttt{contactPoint}, \texttt{normal} и \texttt{penetrationDepth} позволило перейти от простого отталкивания к более реалистичному импульсному отклику.

\subsection{Интегрирование движения и гравитации}
Физический шаг начинается с интегрирования динамических тел. Если для тела включён флаг \texttt{useGravity}, к ускорению добавляется постоянный вектор \texttt{gravity}. После этого обновляется линейная и угловая скорость, применяется демпфирование, ограничение максимальной скорости и изменение трансформации.

\begin{lstlisting}[language=C++,caption={Интегрирование тела с учётом гравитации и демпфирования}]
void integrateBody(
    Transform& transform,
    Rigidbody& rigidbody,
    const Collider* collider,
    const PhysicsSystem::Settings& settings,
    const float dtSeconds) {
    if (rigidbody.isStatic || rigidbody.isKinematic || rigidbody.inverseMass <= 0.0F) {
        rigidbody.accumulatedForce = glm::vec3(0.0F);
        rigidbody.accumulatedTorque = glm::vec3(0.0F);
        rigidbody.velocity = glm::vec3(0.0F);
        rigidbody.angularVelocity = glm::vec3(0.0F);
        return;
    }

    glm::vec3 totalAcceleration = rigidbody.acceleration;
    if (rigidbody.useGravity) {
        totalAcceleration += settings.gravity;
    }
    totalAcceleration += rigidbody.accumulatedForce * rigidbody.inverseMass;

    rigidbody.velocity += totalAcceleration * dtSeconds;
    rigidbody.velocity *= 1.0F / (1.0F + rigidbody.linearDamping * dtSeconds);

    const glm::mat3 inverseInertia =
        worldInverseInertiaTensor(transform, collider, rigidbody);
    rigidbody.angularVelocity += inverseInertia * rigidbody.accumulatedTorque * dtSeconds;
    rigidbody.angularVelocity *= 1.0F / (1.0F + rigidbody.angularDamping * dtSeconds);

    transform.position += rigidbody.velocity * dtSeconds;
    transform.rotationEulerRadians += rigidbody.angularVelocity * dtSeconds;
    rigidbody.accumulatedForce = glm::vec3(0.0F);
    rigidbody.accumulatedTorque = glm::vec3(0.0F);
}
\end{lstlisting}

По сравнению с минимальной схемой из методички здесь дополнительно учитываются угловая динамика, мировая матрица инерции и стабилизация численно неустойчивых значений.

\subsection{Разрешение столкновений}
После построения потенциальных пар \texttt{PhysicsSystem} выполняет несколько solver-итераций. Для каждого контакта сначала проводится мягкая position correction, затем вычисляется нормальный импульс с учётом массы и коэффициента восстановления, а после этого --- тангенциальный импульс трения.

\begin{lstlisting}[language=C++,caption={Разрешение столкновения: correction + impulse + friction}]
void resolveCollision(BodyProxy& bodyA, BodyProxy& bodyB, const CollisionContact& contact) {
    constexpr float kPenetrationSlop = 0.01F;
    constexpr float kCorrectionPercent = 0.34F;
    constexpr float kBounceSpeedThreshold = 1.15F;

    const float inverseMassA = effectiveInverseMass(bodyA.rigidbody);
    const float inverseMassB = effectiveInverseMass(bodyB.rigidbody);
    const float inverseMassSum = inverseMassA + inverseMassB;

    if (inverseMassSum > 0.0F) {
        const float correctionDepth =
            std::max(contact.penetrationDepth - kPenetrationSlop, 0.0F);
        const glm::vec3 correction = contact.normal *
                                     (correctionDepth * kCorrectionPercent);
        if (bodyA.transform != nullptr && inverseMassA > 0.0F) {
            bodyA.transform->position -= correction * (inverseMassA / inverseMassSum);
        }
        if (bodyB.transform != nullptr && inverseMassB > 0.0F) {
            bodyB.transform->position += correction * (inverseMassB / inverseMassSum);
        }
    }

    const glm::vec3 relativeVelocity =
        velocityAtPoint(bodyB, contact.contactPoint) -
        velocityAtPoint(bodyA, contact.contactPoint);
    const float closingSpeed = glm::dot(relativeVelocity, contact.normal);
    if (closingSpeed >= 0.0F) {
        return;
    }

    float restitution = 0.5F * (bodyA.rigidbody->restitution + bodyB.rigidbody->restitution);
    if (-closingSpeed < kBounceSpeedThreshold) {
        restitution = 0.0F;
    }
    ...
}
\end{lstlisting}

Это решение выглядит сложнее, чем ``простейшее отталкивание'' из задания, но даёт заметно более правдоподобное поведение объектов: различающуюся реакцию лёгких и тяжёлых тел, вращение при смещённом ударе и катание сфер по поверхности.

\subsection{Broadphase, sleeping и публикация событий}
Для предотвращения полного перебора всех пар на сцене используется простое пространственное разбиение по трёхмерной сетке (\texttt{broadphaseCellSize}). После расчёта коллизий система сравнивает текущий набор контактов с предыдущим кадром и публикует события \texttt{CollisionEnter}, \texttt{CollisionStay} и \texttt{CollisionExit} в общий \texttt{EventBus}.

\begin{lstlisting}[language=C++,caption={Обновление \code{PhysicsSystem}: broadphase, solver, события}]
void PhysicsSystem::update(World& world, const double dt, core::EventBus& eventBus) {
    world.forEach<Transform, Rigidbody>([&](EntityId entity, Transform& transform, Rigidbody& rigidbody) {
        integrateBody(transform, rigidbody, world.getComponent<Collider>(entity), m_settings, dtSeconds);
    });

    std::vector<BodyProxy> bodies;
    world.forEach<Transform, Collider>([&](EntityId entity, Transform& transform, Collider& collider) {
        BodyProxy proxy{};
        proxy.entity = entity;
        proxy.transform = &transform;
        proxy.rigidbody = world.getComponent<Rigidbody>(entity);
        proxy.collider = &collider;
        if (!computeWorldColliderShape(world, entity, collider, proxy.worldShape)) {
            return;
        }
        proxy.worldBounds = computeWorldBounds(proxy.worldShape);
        bodies.push_back(proxy);
    });

    const std::vector<PotentialPair> potentialPairs =
        buildBroadphasePairs(bodies, m_settings.broadphaseCellSize);

    for (std::size_t iteration = 0; iteration < m_settings.solverIterations; ++iteration) {
        for (const PotentialPair& potentialPair : potentialPairs) {
            CollisionContact contact{};
            if (!intersectColliders(..., contact)) {
                continue;
            }
            resolveCollision(bodyA, bodyB, contact);
        }
    }

    for (const auto& [pair, contact] : frameContacts) {
        if (m_activeContacts.find(pair) == m_activeContacts.end()) {
            eventBus.publish(CollisionEnterEvent{contact});
        } else {
            eventBus.publish(CollisionStayEvent{contact});
        }
    }
}
\end{lstlisting}

В сравнении с минимальным вариантом задания добавление broadphase и sleeping позволило значительно уменьшить просадки FPS при массовом спавне объектов в демонстрационной сцене.

\subsection{Система событий столкновений}
Событийная подсистема реализована как лёгкий \texttt{EventBus} поверх \texttt{std::function} и \texttt{type\_index}. Такой подход хорошо подходит учебному движку: он не требует сложной иерархии сообщений, но даёт типобезопасную подписку и публикацию.

\begin{lstlisting}[language=C++,caption={Структуры событий и шаблонный \code{EventBus}}]
struct CollisionEnterEvent final { CollisionContact contact; };
struct CollisionStayEvent final  { CollisionContact contact; };
struct CollisionExitEvent final  { CollisionContact contact; };

template <class T, class Func>
void EventBus::subscribe(Func&& handler) {
    auto typedHandler = std::function<void(const T&)>(std::forward<Func>(handler));
    m_handlers[typeid(T)].emplace_back(
        [wrapped = std::move(typedHandler)](const void* rawEvent) {
            wrapped(*static_cast<const T*>(rawEvent));
        });
}

template <class T>
void EventBus::publish(const T& event) const {
    const auto it = m_handlers.find(typeid(T));
    if (it == m_handlers.end()) {
        return;
    }
    for (const auto& handler : it->second) {
        handler(&event);
    }
}
\end{lstlisting}

В \texttt{GameplayState} эти события используются для обновления счётчиков столкновений и вывода отладочного текста в ImGui-панель.

\subsection{InputManager и управление камерой}
Ввод вынесен в \texttt{platform::InputManager}. Он хранит текущее и предыдущее состояние клавиш/кнопок, дельту мыши, прокрутку колеса и флаги захвата ввода интерфейсом. Дополнительно предусмотрены ``raw'' методы, игнорирующие ImGui-capture, что было важно для возврата управления камерой без лишнего клика по окну.

\begin{lstlisting}[language=C++,caption={Обновление \code{InputManager} по событиям окна}]
void InputManager::onEvent(const Event& event) {
    switch (event.type) {
        case EventType::KeyPressed:
            s_currentKeyStates[static_cast<std::size_t>(event.key)] = true;
            break;
        case EventType::KeyReleased:
            s_currentKeyStates[static_cast<std::size_t>(event.key)] = false;
            break;
        case EventType::MouseButtonPressed:
            s_currentMouseStates[static_cast<std::size_t>(event.mouseButton)] = true;
            break;
        case EventType::MouseButtonReleased:
            s_currentMouseStates[static_cast<std::size_t>(event.mouseButton)] = false;
            break;
        case EventType::MouseMoved:
            s_mouseDeltaX += event.mouseX - s_mouseX;
            s_mouseDeltaY += event.mouseY - s_mouseY;
            s_mouseX = event.mouseX;
            s_mouseY = event.mouseY;
            break;
        case EventType::MouseScrolled:
            s_scrollDeltaX += event.scrollX;
            s_scrollDeltaY += event.scrollY;
            break;
        default:
            break;
    }
}
\end{lstlisting}

Поверх \texttt{InputManager} реализован \texttt{FreeCameraController}, работающий в стиле fly-camera: мышь вращает камеру при удержании правой кнопки, клавиши \texttt{WASDQE} двигают её в пространстве, а \texttt{Shift} ускоряет полёт.

\begin{lstlisting}[language=C++,caption={Обновление свободной камеры}]
void FreeCameraController::update(
    Transform& transform,
    double dt,
    bool flyModeActive,
    bool mouseLookActive) {
    if (mouseLookActive) {
        const auto [deltaX, deltaY] = platform::InputManager::mouseDeltaRaw();
        transform.rotationEulerRadians.y += static_cast<float>(deltaX) * m_settings.mouseSensitivity;
        transform.rotationEulerRadians.x -= static_cast<float>(deltaY) * m_settings.mouseSensitivity;
        transform.rotationEulerRadians.x = std::clamp(
            transform.rotationEulerRadians.x,
            -m_settings.pitchLimitRadians,
            m_settings.pitchLimitRadians);
    }

    const float yaw = transform.rotationEulerRadians.y;
    const glm::vec3 forward(std::sin(yaw), 0.0F, std::cos(yaw));
    const glm::vec3 worldUp(0.0F, 1.0F, 0.0F);
    const glm::vec3 right = glm::normalize(glm::cross(forward, worldUp));

    glm::vec3 movement(0.0F);
    if (platform::InputManager::isKeyDownRaw(platform::KeyCode::W)) movement += forward;
    if (platform::InputManager::isKeyDownRaw(platform::KeyCode::S)) movement -= forward;
    if (platform::InputManager::isKeyDownRaw(platform::KeyCode::A)) movement -= right;
    if (platform::InputManager::isKeyDownRaw(platform::KeyCode::D)) movement += right;
    if (platform::InputManager::isKeyDownRaw(platform::KeyCode::E)) movement += worldUp;
    if (platform::InputManager::isKeyDownRaw(platform::KeyCode::Q)) movement -= worldUp;
    ...
}
\end{lstlisting}

Особый момент интеграции с интерфейсом: в \texttt{Application::onEvent()} добавлен \texttt{gameplayLookOverride}, благодаря которому удержание правой кнопки мыши снова передаёт управление gameplay даже после взаимодействия с ImGui.

\subsection{Debug-визуализация коллайдеров}
Для AABB используется отрисовка каркаса через \texttt{Renderer::drawDebugAabb()}, а для сфер --- три кольца по осям. Это решение оказалось визуально чище, чем полупрозрачная заливка или второй solid mesh поверх сферы.

\begin{lstlisting}[language=C++,caption={DebugRenderSystem для AABB и сферических коллайдеров}]
void DebugRenderSystem::render(
    World& world,
    renderer::RenderAdapter& renderer,
    const glm::mat4& viewProjectionMatrix) {
    world.forEach<Transform, Collider>([&](EntityId entity, Transform& transform, Collider& collider) {
        WorldColliderShape worldShape{};
        if (!computeWorldColliderShape(world, entity, collider, worldShape)) {
            return;
        }

        const WorldAabb bounds = computeWorldBounds(worldShape);
        const auto* rigidbody = world.getComponent<Rigidbody>(entity);
        const glm::vec4 color =
            (rigidbody != nullptr && rigidbody->isDynamic())
                ? glm::vec4(0.25F, 0.95F, 0.45F, 1.0F)
                : glm::vec4(0.95F, 0.7F, 0.2F, 1.0F);

        if (worldShape.type == ColliderType::Sphere) {
            drawDebugSphereRings(renderer, worldShape.sphere, viewProjectionMatrix, color);
            return;
        }

        renderer.drawDebugAabb(bounds.min, bounds.max, viewProjectionMatrix, color);
    });
}
\end{lstlisting}

В актуальной версии движка переключение режима отладочной визуализации привязано к клавише \texttt{F1}, а не к \texttt{F3}. Это решение отражено и в ImGui-панели, и в логике \texttt{GameplayState}.

\subsection{Демонстрационная сцена и отладочный UI}
Для проверки работы всех подсистем создана отдельная crash-sandbox сцена. Она включает пол, стены, колонны, ``ворота'' из кубов, ряд домино, заднюю пирамидальную группу, тяжёлый страйкер и отдельную сферу. Через ImGui можно запускать симуляцию, перезапускать сцену, спавнить каскад объектов, менять гравитацию, число solver-итераций и размер broadphase-сетки.

\begin{lstlisting}[language=C++,caption={Ключевые элементы \code{GameplayState::renderUi()}}]
ImGui::TextUnformatted("Interactive crash sandbox");
ImGui::TextUnformatted("Hold RMB + WASDQE for UE-style fly camera");
ImGui::TextUnformatted("Mouse wheel changes speed, Shift boosts movement, Space launches striker");
ImGui::TextUnformatted("F1 toggles collider debug. Scene stays posed until physics starts.");

if (!m_simulationRunning) {
    if (ImGui::Button("Start Physics")) {
        startSimulation();
    }
} else if (ImGui::Button("Pause Physics")) {
    pauseSimulation();
}
if (ImGui::Button("Reset Scene")) {
    resetDemoScene();
}
if (ImGui::Button("Launch Striker")) {
    launchStriker();
}
if (ImGui::Button("Drop Wave")) {
    spawnDropWave();
}

ImGui::Text("Bodies / sleeping / broadphase pairs: %zu / %zu / %zu",
    m_physicsSystem.bodyCount(),
    m_physicsSystem.sleepingBodyCount(),
    m_physicsSystem.broadphasePairCount());
\end{lstlisting}

Именно эта сцена использовалась как основной стенд для экспериментальной проверки устойчивости физики, качества управления камерой и визуальной отладки.

\subsection{Автоматические тесты}
Помимо ручной проверки в приложении, для физики добавлен отдельный тестовый таргет \texttt{engine\_physics\_tests}. Он проверяет посадку коробки на пол, зависимость разлёта от массы, появление угловой скорости при смещённом ударе и корректное качение сферы.

\begin{lstlisting}[language=C++,caption={Тест качения сферы по полу}]
bool runSphereRollingTest() {
    engine::ecs::World world;
    engine::core::EventBus eventBus;
    engine::ecs::PhysicsSystem physics;

    const auto floor = spawnBox(world, glm::vec3(0.0F, -0.5F, 0.0F),
                                glm::vec3(14.0F, 1.0F, 14.0F), true, 1.0F);
    const auto sphere = spawnSphere(world, glm::vec3(-2.0F, 0.52F, 0.0F),
                                    1.0F, false, 1.2F);
    (void)floor;

    auto* sphereBody = world.getComponent<engine::ecs::Rigidbody>(sphere);
    auto* sphereTransform = world.getComponent<engine::ecs::Transform>(sphere);
    sphereBody->velocity = glm::vec3(4.0F, 0.0F, 0.0F);

    for (int step = 0; step < 180; ++step) {
        physics.update(world, 1.0 / 60.0, eventBus);
    }

    return sphereTransform->position.x > 1.0F &&
           std::abs(sphereTransform->position.y - 0.5F) < 0.12F &&
           std::abs(sphereBody->angularVelocity.z) > 0.2F;
}
\end{lstlisting}

\section{Результаты}
\subsection{Функциональный результат}
В результате выполнения работы движок демонстрирует:
\begin{itemize}
  \item работу компонентов \texttt{Rigidbody} и \texttt{Collider} внутри ECS;
  \item применение гравитации и интегрирование линейного/углового движения;
  \item обнаружение столкновений для AABB и сфер;
  \item разрешение коллизий с учётом массы, восстановления и трения;
  \item публикацию событий \texttt{CollisionEnter}, \texttt{CollisionStay}, \texttt{CollisionExit};
  \item свободную 3D-камеру с управлением мышью и клавиатурой;
  \item debug-отрисовку коллайдеров;
  \item оптимизацию столкновений через broadphase grid и sleeping;
  \item отдельные автоматические тесты физической подсистемы.
\end{itemize}

\subsection{Сопоставление с требованиями задания}
\begin{table}[H]
\centering
\begin{tabular}{p{0.47\linewidth} p{0.43\linewidth}}
\toprule
\textbf{Требование задания} & \textbf{Реализация в проекте} \\
\midrule
\texttt{Rigidbody} и \texttt{Collider} & Реализованы, плюс угловая динамика, restitution, friction, sleeping \\
\texttt{PhysicsSystem} с гравитацией и коллизиями & Реализована, расширена до импульсной модели и solver-итераций \\
\texttt{InputManager} & Реализован в \code{engine/platform/Input.*}, поддерживает raw/captured input \\
События столкновений & Реализованы через \texttt{EventBus} и события Enter/Stay/Exit \\
Debug-отрисовка коллизий & Реализована для AABB и сферических коллайдеров \\
Управление 3D-камерой & Реализовано через \texttt{FreeCameraController} и \texttt{GameplayState} \\
Демо-сцена & Реализована crash-sandbox сцена с ImGui-параметрами и массовым спавном \\
\bottomrule
\end{tabular}
\caption{Соответствие требованиям практического задания}
\label{tab:requirements}
\end{table}

\subsection{Иллюстрации и проверка работоспособности}
В репозитории отчёта не хранится фиксированный набор PNG-скриншотов для ПЗ4, поэтому в текстовой версии отчёта основной упор сделан на архитектуру, листинги и результаты автоматических тестов. При необходимости к печатной или финальной электронной версии можно добавить два актуальных скриншота:
\begin{enumerate}
  \item сцена в обычном режиме рендера;
  \item сцена в режиме debug-отрисовки коллайдеров.
\end{enumerate}

С практической точки зрения работоспособность подтверждается следующими наблюдениями:
\begin{itemize}
  \item динамические тела падают на пол и больше не проникают сквозь него;
  \item лёгкие объекты разлетаются сильнее тяжёлых;
  \item смещённый удар создаёт угловую скорость;
  \item сферы поддерживаются полом и катятся, накапливая вращение;
  \item F1 корректно переключает визуализацию коллайдеров;
  \item RMB + WASDQE обеспечивает стабильное управление камерой без дополнительного клика после ImGui.
\end{itemize}

\section{Возникшие трудности и способы их решения}
\begin{itemize}
  \item \textbf{Нестабильность простого отталкивания.} Базовый positional push выглядел неестественно и вызывал сильные разлёты. Решение --- добавить импульсную модель, трение, коэффициент восстановления и несколько solver-итераций.

  \item \textbf{Просадки производительности при массовом спавне.} Полный перебор пар коллизий быстро становился узким местом. Решение --- broadphase grid, sleeping и бюджетированная очередь спавна в \texttt{GameplayState}.

  \item \textbf{Конфликт ввода gameplay и ImGui.} После взаимодействия с UI приходилось повторно кликать по окну. Решение --- raw input path и \texttt{gameplayLookOverride} при удержании правой кнопки мыши.

  \item \textbf{Наглядная визуализация сферических коллайдеров.} Сплошная отладочная сфера давала артефакты наложения. Решение --- отрисовка трёх тонких колец по осям.

  \item \textbf{Корректное отображение катящейся сферы.} Ранний OBJ-меш сферы имел дефекты на полюсах. Решение --- переход на процедурную сферу в рендерере и синхронизация визуальной формы с коллайдером.
\end{itemize}

\section{Ответы на контрольные вопросы}
\begin{enumerate}
  \item \textbf{В чём разница между \texttt{OnCollision} и \texttt{OnTrigger} в контексте физики?}\\
  \texttt{OnCollision} срабатывает при реальном физическом контакте, который влияет на позиции и скорости тел. \texttt{OnTrigger} обычно вызывается для ``проходного'' объёма: пересечение фиксируется, но импульсного разрешения и физического отталкивания не происходит.

  \item \textbf{Почему для физики предпочтительнее фиксированный шаг обновления?}\\
  Фиксированный шаг делает поведение симуляции стабильнее и предсказуемее. При переменном \textit{dt} скачки времени кадра могут приводить к нестабильным коллизиям, разной высоте прыжка и неповторяемым результатам симуляции.

  \item \textbf{Как рассчитать силу гравитации, чтобы объекты падали с реалистичным ускорением?}\\
  В упрощённой модели достаточно использовать постоянный вектор ускорения \(g = (0, -9.81, 0)\) м/с\(^2\) и каждый шаг обновлять скорость по формуле \(v = v + g \cdot dt\). Если нужна именно сила, то \(F = m \cdot g\).

  \item \textbf{Какие существуют стратегии разрешения коллизий (\textit{position correction}, \textit{impulse-based})?}\\
  \textit{Position correction} просто разводит пересёкшиеся тела по нормали контакта и удобна как минимальное решение. \textit{Impulse-based} дополнительно меняет скорости и вращение тел на основе массы, нормали контакта, восстановления и трения, поэтому даёт более правдоподобный результат.

  \item \textbf{Как организовать маппинг действий (например, ``прыжок'') на разные клавиши?}\\
  Обычно ввод делится на два слоя: низкоуровневый опрос устройств и уровень действий. На втором уровне действие \texttt{Jump} сопоставляется списку клавиш/кнопок, а геймплей код проверяет не конкретную клавишу, а действие. В текущем учебном проекте для простоты используется прямой опрос клавиш, но \texttt{InputManager} уже является хорошей базой для такого расширения.

  \item \textbf{Для чего нужно пространственное разбиение при обнаружении коллизий?}\\
  Оно уменьшает число потенциальных пар. Вместо проверки всех \(n(n-1)/2\) комбинаций движок сначала группирует объекты по ячейкам, а затем тестирует только локально близкие тела. Это критично при росте числа динамических объектов.

  \item \textbf{Как преобразовать локальный AABB коллайдера в мировое пространство для отрисовки?}\\
  Нужно взять мировую матрицу сущности, трансформировать центр коллайдера с учётом \texttt{offset}, затем вычислить мировые half-extents через абсолютную матрицу поворота-масштаба. После этого можно получить \texttt{min = center - extents} и \texttt{max = center + extents}.

  \item \textbf{Какие недостатки имеет AABB по сравнению с OBB (Oriented Bounding Box)?}\\
  AABB проста и быстра, но плохо описывает повернутые объекты: при вращении объём становится слишком большим и даёт ложные срабатывания. OBB точнее аппроксимирует ориентированные тела, но требует более сложной математики и дороже в вычислении.
\end{enumerate}

\section{Заключение}
В ходе выполнения практической работы в движок была интегрирована полноценная подсистема базовой физики и современный слой обработки пользовательского ввода. На уровне ECS добавлены компоненты \texttt{Rigidbody}, \texttt{Collider} и \texttt{Camera}; реализована \texttt{PhysicsSystem}, выполняющая гравитацию, интегрирование движения, детектирование столкновений и их разрешение; создана визуальная отладка коллайдеров и свободная 3D-камера. Дополнительно проект был расширен поддержкой сферических коллайдеров, импульсной моделью отклика, broadphase-оптимизацией, sleeping тел и набором автоматических тестов. Полученная архитектура служит хорошей основой для следующих шагов развития движка: триггеров, более точных коллайдеров, constraint-системы и асинхронной физической симуляции.

\subsection*{Репозиторий проекта на GitHub}
\href{https://github.com/TurbikXD/TurboGameEngine}{https://github.com/TurbikXD/TurboGameEngine}.

\end{document}
